\documentclass[11pt,a4paper]{article}

\usepackage[utf8]{inputenc}

\usepackage[T1]{fontenc}

\usepackage[polish]{babel}

\usepackage{lmodern}

\usepackage{geometry}

\geometry{margin=27mm}

\usepackage[hidelinks]{hyperref}

\title{Struktura czynnikowa i trafność narzędzia DERA-HRI do oceny doświadczenia interakcji z humanoidem}

\author{}

\date{Szkic naukowy, wersja 0.1.0}

\begin{document}

\maketitle

\begin{abstract}

Tematem pracy jest Struktura czynnikowa i trafność narzędzia DERA-HRI do oceny doświadczenia interakcji z humanoidem. Tekst systematyzuje uzasadnienie, hipotezy i sposób weryfikacji zagadnienia. Proponowany wskaźnik: Replikowalność modelu czynnikowego DERA-HRI na niezależnej próbie oraz przedziały dla rzetelności i trafności. Nie przedstawia się niezweryfikowanych wyników ani fikcyjnych cytowań.

\end{abstract}

\section{Wprowadzenie}

Wnioskowanie na podstawie pomiaru wymaga walidacji definicji sygnału, kalibracji i kontroli błędów. Zagadnienie Struktura czynnikowa i trafność narzędzia DERA-HRI do oceny doświadczenia interakcji z humanoidem wymaga odróżnienia badanego konstruktu od jego operacjonalizacji i kontekstowych źródeł zmienności. Trzeba porównać możliwe wyjaśnienia i ustalić granice interpretacji dowodów.

\section{Cel i pytania badawcze}

Celem pracy jest rozstrzygnięcie zagadnienia: Struktura czynnikowa i trafność narzędzia DERA-HRI do oceny doświadczenia interakcji z humanoidem. Proponowany wskaźnik: Replikowalność modelu czynnikowego DERA-HRI na niezależnej próbie oraz przedziały dla rzetelności i trafności. Hipotezy i ich kierunki muszą zostać określone przed analizą danych lub krytyczną syntezą źródeł.

\section{Materiał i metody}

Opisać wersje narzędzi, jednostki, źródła znaczników czasu i kryteria jakości. Przed analizą należy określić jednostkę obserwacji, warunki kontrolne, kryteria wyłączenia oraz sposób ograniczania zmiennych zakłócających.  Przed rozpoczęciem badań należy udokumentować warunki etyczne, zgodę, prawa do danych i rzeczywistą dostępność narzędzi.

\section{Plan analizy}

Ocenić powtarzalność, trafność, rozkład błędów i wpływ zmian warunków. Kryteria wyboru modelu analitycznego, analiza niepewności i postępowanie z brakami danych wymagają uprzedniej specyfikacji.

\section{Wyniki}

W niniejszym pliku nie umieszczono wyników empirycznych, statystyk ani nieprzeprowadzonych analiz. Sekcja wymaga uzupełnienia po uzyskaniu i weryfikacji materiału dowodowego.

\section{Dyskusja i ograniczenia}

Ograniczenia obejmują możliwe błędy pomiaru, dobór próby, alternatywne mechanizmy i ograniczoną możliwość uogólniania. Korelacja nie dowodzi przyczynowości; poprawność techniczna nie jest dowodem znaczenia klinicznego.

\section{Wnioski}

Zdefiniowano ramy opracowania, ale nie sformułowano wniosków wykraczających poza dostępny stan dowodów.

\section*{Bibliografia}

Weryfikacja literatury i przypisanie cytowań do stwierdzeń są kolejnym etapem prac.

\end{document}

